\documentclass[a4paper]{ltjsarticle}
\usepackage{luatexja-fontspec}
\usepackage{amsmath, amssymb, amsthm}
\usepackage{tikz-cd}
\usepackage{hyperref}

\title{P1\_Extended\_Next の数学的背景解説}
\author{CODEX (OpenAI)\\TeX ファイル生成および Lean ファイル生成担当}
\date{\today}

\begin{document}

\maketitle

\begin{abstract}
本稿は \texttt{P1\_Extended\_Next.lean} の主要概念を学部生向けに解説したメモである。
ブルバキ流の構造主義的視点を踏まえ、環論・ブール代数・フィルター・測度論の各節で
どのような ``塔 (StructureTower)'' が構成されているかを紹介する。
\end{abstract}

\section{導入}

Lean 4 形式化ファイル \texttt{P1\_Extended\_Next.lean} では、以下のテーマが統一的に扱われている。
\begin{itemize}
  \item 環準同型の像とイデアル格子
  \item ブール代数の補元構造と Stone 双対性への動機付け
  \item フィルター・フィルトレーションを StructureTower で捉える視点
  \item 素朴な $\sigma$-代数から測度論的塔を得る枠組み
\end{itemize}
本メモはそれぞれの節を学部レベルの数式と図式で要約する。

\section{環とイデアルの塔}

Lean ファイルでは、環準同型 $f:R \to S$ の像が部分環になることを
\[
  \mathrm{imageSubring}(f) := f(R) \subseteq S
\]
として \texttt{Subring} に格納している。さらに、イデアル全体を前半順序付き集合とみなして
StructureTower を構成し、各レベルがイデアル $I \lhd R$ に対応する。

\subsection*{第一同型定理の位置づけ}
環の第一同型定理 $R/\ker f \cong \operatorname{Im} f$ を念頭に置くと、
以下のような図式が自然に現れる。
\[
  \ker f \subseteq R \xrightarrow{f} S,\qquad R/\ker f \xrightarrow{\cong} \operatorname{Im} f.
\]
Lean 側では \texttt{RingHom.quotientKerEquivRange f} がこの同型を実装している。

\subsection*{イデアル塔の例}
StructureTower を使うと、各元 $x\in R$ が自身の生成するイデアル $\langle x\rangle$ のレベルに属すること、
全レベルの合併が $R$ 全体を覆うことが簡潔に示せる。これは「元 $\mapsto$ 包含するイデアル」という
ブルバキ風の ``母構造'' を視覚化したものである。

\section{ブール代数と補元不動点}

ブール代数 $B$ に対して、補元演算が自分自身を保つ元の集合
\[
  \mathrm{ComplementFixed}(B) = \{x \in B \mid x^c = x\}
\]
を考えると、$B$ が非自明 ($\bot \neq \top$) である限り空集合となることが Lean で証明されている。
これは $x = x^c$ ならば $x = \bot = \top$ を意味し矛盾するという初等的議論である。

この考察は Stone 双対性に向けた第一歩であり、補元がどのように
位相的な ``clopen 集合'' の振る舞いと平行するかを示している。

\section{フィルターとフィルトレーション}

集合 $X$ 上のフィルター $F$ を用いて、レベルを
\[
  \mathrm{level}(A) = \{x \in X \mid A \in F,\ x \in A\}
\]
と定義すると StructureTower を得る。ここで $A$ はフィルターに属する集合である。
これは「大きい集合」を段階的に眺める際の便利な抽象化である。

\subsection*{有限フィルトレーションの例}
$\mathbb{N}$ の場合、$F_n = \{k \in \mathbb{N} \mid k \le n\}$ とおくと
単調性 $F_n \subseteq F_{n+1}$ が成り立ち、限界において全体 $\mathbb{N}$ を覆う。
Lean ファイルでも \texttt{natFiniteFiltration} として同様の構成が導入されている。

\section{測度論的構造への足がかり}

最後の節では、集合族 $g \subseteq \mathcal{P}(\alpha)$ から生成される $\sigma$-代数
$\sigma(g)$ を文字通りの定義に従って構成し、これに自然な順序を入れて StructureTower を得る。
さらに、可測空間を既に持っている場合には
\[
  \mathrm{sigmaAlgebraOfMeasurable}(\alpha) = \{A \subseteq \alpha \mid A \text{ が可測}\}
\]
として ``測度論的塔'' を構成している。

\subsection*{図式: 生成された $\sigma$-代数}
以下の図式は、生成過程の概念的流れを示したものである。
\begin{center}
\begin{tikzcd}
g \arrow[hookrightarrow]{r} & \mathcal{P}(\alpha) \arrow[d, dashed, "\text{閉包操作}"] \\
 & \sigma(g)
\end{tikzcd}
\end{center}
点線矢印は、補集合と可算和に関する閉包を施して得られる最小の $\sigma$-代数を表す。

\section{まとめと今後の見通し}

本稿では Lean 実装の骨格を学部生向けに解説した。今後は以下の発展が期待される。
\begin{itemize}
  \item Stone 双対性の完全な形式化（Ultrafilter Stone 空間と clopen アルジェブラの同型）
  \item 連続実数値函数の前層が Lean 上でもシーフ条件を満たすことの証明
  \item StructureTower をカテゴリ的視点から解析し、普遍性を明示する
\end{itemize}
これらが実現されれば、ブルバキ的 ``母構造'' をさらに高次元で探索する土台が整う。

\end{document}
